\section{课后练习}

\subsection{Packaging and Shipping Code}

\subsubsection{用 pip 安装与卸载第三方包}
pip 是 Python 的包管理工具，可在线安装、卸载并查看已安装的包，相关内容可参考课程资源中的 runoob Python 教程。
\begin{minted}[frame=single, fontsize=\small]{bash}
pip install requests          # 安装包
pip install requests==2.31.0  # 安装指定版本
pip uninstall requests -y     # 卸载包
pip list                      # 查看已安装的包
\end{minted}

\subsubsection{用 python -m build 构建 wheel 与 sdist}
python -m build 会依据 pyproject.toml 同时生成 wheel 和源码包 sdist，是发布 Python 包的标准方式。
\begin{minted}[frame=single, fontsize=\small]{bash}
python3 -m build
# 生成 dist/*.whl 与 dist/*.tar.gz
\end{minted}

\subsubsection{从本地 wheel 安装包}
在离线或已构建场景下，可直接用 pip 从本地 .whl 文件安装，而不依赖网络源。
\begin{minted}[frame=single, fontsize=\small]{bash}
pip install dist/greetlab-25020007136-0.1.0-py3-none-any.whl
\end{minted}

\subsection{智能体编程}

\subsubsection{测试驱动的智能体修复流程}
先编写一个失败测试明确期望行为，再让智能体以该测试为验收标准修改代码，可形成可验证的修复闭环。
\begin{minted}[frame=single, fontsize=\small]{bash}
# 1. 先写失败测试并运行，确认失败
pytest test_cli.py
# 2. 让智能体修改实现，以"运行 pytest 通过"为约束
# 3. 再次运行测试，确认修复有效
pytest test_cli.py
\end{minted}

\subsubsection{用 git diff 审查智能体的改动}
智能体可能引入与目标无关的改动，提交前应使用 git diff 检查，撤销无关修改。
\begin{minted}[frame=single, fontsize=\small]{bash}
git diff                # 查看工作区与暂存区的差异
git diff --stat         # 仅查看改动文件与行数
git checkout -- <file>  # 撤销某个文件的无关修改
\end{minted}

\subsection{不止于代码}

\subsubsection{编写规范的 Git 提交信息}
采用 Conventional Commits 约定，标题用祈使语气并标注类型（fix/feat/docs 等），正文说明问题与方案，便于历史回溯。
\begin{minted}[frame=single, fontsize=\small]{bash}
git commit -m "fix: 拦截纯空白姓名的无效输入" \
           -m "当前 --name 未校验空白字符，导致错误问候且退出码为 0。"
\end{minted}

\subsubsection{编写可复现的 Issue 报告}
一份可执行的 Issue 应包含环境、复现命令、期望结果与实际结果，未知信息明确写“待确认”，便于他人快速定位。
\begin{minted}[frame=single, fontsize=\small]{markdown}
### Issue: 空姓名仍输出问候语
- 环境：Windows 10, Python 3.10+
- 复现命令：sdt-greet --name "   "
- 预期结果：报错并以非零状态码退出。
- 实际结果：输出问候语并以 0 退出。
\end{minted}

\subsection{Python 与 PyTorch}

\subsubsection{PyTorch 张量的创建与运算}
张量是 PyTorch 的基本数据结构，支持与 NumPy 类似的创建方式与广播运算，入门可参考课程资源中的 PyTorch 官方教程。
\begin{minted}[frame=single, fontsize=\small]{python}
import torch
a = torch.linspace(-1, 1, 5)   # 等间隔张量
x = a.reshape(-1, 1)           # 改变形状
y = 3 * x - 1                  # 广播运算
print(y)
\end{minted}

\subsubsection{自动求导与梯度清零}
PyTorch 通过 backward 自动计算梯度，但梯度默认累加，因此每次迭代前需调用 zero\_grad 清空。
\begin{minted}[frame=single, fontsize=\small]{python}
import torch
x = torch.tensor([2.0], requires_grad=True)
y = x * x                       # y = x^2
y.backward()                    # 自动求导 dy/dx = 2x
print(x.grad)                   # tensor([4.])
opt.zero_grad()                 # 清空累积的梯度
\end{minted}

\subsubsection{用 nn.Linear 搭建最小训练循环}
nn.Linear 表示线性层，配合损失函数与优化器即可构成最小训练循环，三步完成一次参数更新。
\begin{minted}[frame=single, fontsize=\small]{python}
import torch
from torch import nn

model = nn.Linear(1, 1)
loss_fn = nn.MSELoss()
opt = torch.optim.SGD(model.parameters(), lr=0.1)

pred = model(x)
loss = loss_fn(pred, y)
opt.zero_grad()   # 清空梯度
loss.backward()   # 反向传播
opt.step()        # 更新参数
\end{minted}